% SPDX-License-Identifier: CC-BY-4.0
\documentclass[11pt,letterpaper]{article}

\usepackage[T1]{fontenc}
\usepackage{amsmath,amssymb}
\usepackage{tikz}
\usepackage{pgfplots}
\pgfplotsset{compat=1.18}
\usepackage{algorithm}
\usepackage{algpseudocode}
\usepackage{booktabs}
\usepackage{enumitem}
\usepackage[hidelinks]{hyperref}

\newcommand{\N}{\mathbb{N}}
\newcommand{\Z}{\mathbb{Z}}
\newcommand{\Q}{\mathbb{Q}}
\newcommand{\R}{\mathbb{R}}
\newcommand{\C}{\mathbb{C}}
\newcommand{\Map}[3]{#1 \colon #2 \to #3}
\newcommand{\Mod}[1]{#1\text{-}\mathrm{Mod}}
\newcommand{\GroupHom}{\mathrm{Hom}_{\mathbf{Grp}}}
\newcommand{\RingHom}{\mathrm{Hom}_{\mathbf{Ring}}}

\newtheorem{theorem}{Theorem}
\newtheorem{lemma}[theorem]{Lemma}
\newtheorem{definition}[theorem]{Definition}
\newtheorem{remark}[theorem]{Remark}
\newenvironment{proof}{\noindent\textit{Proof.}}{\hfill\(\square\)}

\title{Generalized Collatz Library}
\author{ascendnoosphere}
\date{\today}

\begin{document}
\maketitle

\section{Signed raw dynamics and legal words}

% The declared domain matches Parameters and does not silently assume positivity.
\begin{definition}
Parameters are integers \(q,c,d\) with \(|d|\ge2\). On all of \(\Z\), define
\[
F(n)=\begin{cases}
n/d,&d\mid n,\\
qn+c,&d\nmid n.
\end{cases}
\]
Division is exact on its branch. The sign of \(d\) is retained, and division has
precedence at zero. The two instructions \(A,D\) define partial maps
\[
A(n)=qn+c\quad(d\nmid n),\qquad D(n)=n/d\quad(d\mid n).
\]
Outside the displayed domain the instruction is undefined.
\end{definition}

% The pointwise case proof includes all undefined intermediate states.
\begin{lemma}[Partial composition]
For partial maps, let \(f;g\) mean execute \(f\) and then \(g\): it is undefined when
\(f\) is undefined or when \(g\) is undefined at its output. This composition is
associative, and the everywhere-defined identity is a left and right identity.
\end{lemma}
\begin{proof}
Fix an input \(x\). If \(f(x)\) is undefined, both parenthesizations are undefined.
Otherwise write \(f(x)=y\). If \(g(y)\) is undefined, both are again undefined.
Otherwise write \(g(y)=z\); both parenthesizations are exactly \(h(z)\), including
its possible undefinedness. The identity changes neither the input nor any output,
so it also preserves undefinedness. Pointwise equality proves all three laws.
\end{proof}

% Concatenation keeps a single intermediate witness common to both words.
\begin{theorem}[Guarded concatenation]
For chronological words \(u,v\), executing \(uv\) means executing \(u\) first.
It succeeds from \(n\) to \(m\) exactly when there is an integer \(k\) such that
\(u\) succeeds from \(n\) to \(k\) and \(v\) succeeds from \(k\) to \(m\).
\end{theorem}
\begin{proof}
The empty word is the identity. For a word beginning with one instruction, its
execution is that instruction followed by the tail execution. Induction on the
length of \(u\), using partial associativity, gives
\(F_{uv}=F_u;F_v\). The definition of partial composition then gives exactly the
stated common intermediate witness. No affine relaxation was used.
\end{proof}

% Raw recovery proves the finite-to-word bridge in both directions.
\begin{theorem}[Exact raw itinerary recovery]
If a word \(w\) succeeds from \(n\) to \(m\), then
\(m=F^{|w|}(n)\), and \(w\) is precisely the first \(|w|\) raw branch choices of
that orbit. Conversely, for every \(k\ge0\), those first \(k\) branch choices form
a legal word from \(n\) to \(F^k(n)\).
\end{theorem}
\begin{proof}
At each state, exactly one of \(d\mid n\) and \(d\nmid n\) holds, so exactly one
instruction is legal; its output is \(F(n)\). For a successful nonempty word,
its first instruction is therefore forced. Its tail succeeds from \(F(n)\).
Induction on the tail length gives both its forced itinerary and its endpoint.
The iteration identity \(F^k(F(n))=F^{k+1}(n)\) follows by induction on \(k\).
For the converse, choose the legal instruction at the current state and repeat.
Induction on \(k\) proves legality, length \(k\), and the stated endpoint.
The case \(k=0\) is the empty word at \(n\).
\end{proof}

\section{Affine presentations and the additive offset}

% The denominator is nonzero, but the presentation has no arithmetic guards.
\begin{definition}
An affine presentation is integer data \(a=(P,B,Q)\) with \(Q\ne0\), acting on
\(\Q\) by \(a(x)=(Px+B)/Q\). Here \(B\) is the \emph{full} offset.
The identity presentation is \((1,0,1)\).
\end{definition}

% These unreduced data are the same fields used by AffineThen.
\begin{theorem}[Composition]
For two presentations, execution of the first followed by the second gives
\[
(P_2P_1,\ P_2B_1+B_2Q_1,\ Q_2Q_1).
\]
Its rational action is exactly the composition. Unreduced presentation composition
is associative and has \((1,0,1)\) as a two-sided identity.
\end{theorem}
\begin{proof}
The product denominator is nonzero. Direct substitution gives
\[
\frac{P_2(P_1x+B_1)/Q_1+B_2}{Q_2}
=\frac{P_2P_1x+P_2B_1+B_2Q_1}{Q_2Q_1}.
\]
For three presentations, both parenthesizations have multiplier \(P_3P_2P_1\),
denominator \(Q_3Q_2Q_1\), and offset
\(P_3P_2B_1+P_3B_2Q_1+B_3Q_2Q_1\).
The identity laws follow by substitution in all three fields.
\end{proof}

% The recursion separates the unscaled exponent offset from the full offset.
\begin{definition}
An exponent letter \(a\in\N\) presents \((qx+c)/d^a\), with \(d\ne0\).
For a chronological exponent list \(w\), put \(B_{\varnothing}=0\) and
\[
B_{a::w}=q^{|w|}+d^aB_w,
\]
where \(|w|\) is the length of the tail.
Zero exponents are permitted in this algebraic definition; it is not an
exact accelerated-itinerary definition.
\end{definition}

% The induction also applies to empty, zero-exponent and signed-base cases.
\begin{theorem}[Exponent-word data]
The exponent word \(w\) has presentation
\[
P=q^{|w|},\qquad Q=d^{\sum_{a\in w}a},\qquad B_{\rm full}=cB_w.
\]
Appending lists is chronological affine composition. Its raw expansion---one
affine letter followed by \(a\) division letters for each entry \(a\)---has
the same presentation.
\end{theorem}
\begin{proof}
The empty list gives \((1,0,1)\). Suppose the claim holds for \(w\).
Composing the head letter \((q,c,d^a)\) with the tail gives multiplier
\(q^{|w|}q=q^{|w|+1}\), denominator
\(d^{\sum w}d^a=d^{a+\sum w}\), and full offset
\(q^{|w|}c+cB_wd^a=c(q^{|w|}+d^aB_w)\).
This is the defining recursion. Induction and affine associativity prove the
append law. No positivity, primality or coprimality is needed. Associativity
also shows that the raw expansion presents the same data, so the closure
equation below applies to that expansion.
\end{proof}

% The relaxation recovers values only when the original partial program is legal.
\begin{theorem}[Guarded affine agreement]
Give a raw \(A\)-letter the presentation \((q,c,1)\) and a raw \(D\)-letter
\((1,0,d)\), composing chronologically. If the raw word succeeds from integer
\(n\) to integer \(m\), its rational presentation sends \(n\) to \(m\).
\end{theorem}
\begin{proof}
For a legal affine letter the formula is its integer output. For a legal division,
write \(n=dm\); since \(d\ne0\), rational division gives \(n/d=m\), including
negative \(d\). The empty word agrees with the identity. A legal nonempty word
has one legal first output and a legal tail. Induction and the affine composition
identity prove agreement for the whole word.
\end{proof}

\section{Closure and the scalar conductor}

% Neutral determinants remain in the equation rather than being divided away.
\begin{theorem}[Fixed-point equation]
Let \(\Delta=Q-P\). A rational \(x\) is fixed by \((P,B,Q)\) exactly when
\(\Delta x=B\). If \(\Delta\ne0\), its unique rational fixed point is
\(B/\Delta\). If a raw word is legally closed at integer \(n\), then
\(\Delta n=B\) and hence \(\Delta\mid B\). If that legally closed word is
neutral, \(\Delta=0\), then \(B=0\).
\end{theorem}
\begin{proof}
Since \(Q\ne0\), multiply \((Px+B)/Q=x\) by \(Q\) and rearrange to get
\((Q-P)x=B\); the steps are reversible. A nonzero \(\Delta\) permits division.
For a legally closed raw word, guarded affine agreement supplies the fixed-point
hypothesis. Its integer state is the divisibility witness. At \(\Delta=0\),
the equation is \(0=B\), not an invitation to divide by zero.
\end{proof}

% This criterion is for scalar closure, not for simultaneous intermediate guards.
\begin{theorem}[Scalar conductor criterion]
For integers \(\Delta\ne0,b,c\), let
\[
g=\gcd(|\Delta|,|b|),\qquad \gamma=|\Delta|/g.
\]
Then \(\gamma\mid |c|\) if and only if \(\Delta\mid bc\), if and only if there
exists \(n\in\Z\) with \(\Delta n=cb\).
\end{theorem}
\begin{proof}
The common divisor \(g\) is positive and divides both \(\Delta\) and \(b\).
B\'ezout's identity \cite{Lang02} gives integers \(u,v\) with \(g=\Delta u+bv\).
If \(\Delta\mid bc\), multiplying this identity by \(c\) shows
\(\Delta\mid gc\). Cancelling the nonzero common factor \(g\) gives
\(\Delta/g\mid c\), equivalently \(|\Delta/g|\mid |c|\).
Exact divisibility implies \(|\Delta/g|=|\Delta|/g=\gamma\).
Conversely, \(\Delta/g\mid c\) implies \(\Delta\mid gc\); because \(g\mid b\),
it implies \(\Delta\mid bc\). Finally, integer divisibility is exactly the
existence of the displayed multiplication witness. This proof never uses an
intermediate branch guard, so the conclusion alone cannot certify a cycle.
\end{proof}

\section{Labelled graphs and marked systems}

% Edge identities and labels are not replaced by a simple adjacency relation.
\begin{definition}
A labelled graph is \((V,E,\Sigma,s,t,\ell)\), where \(s,t:E\to V\) and
\(\ell:E\to\Sigma\). A morphism consists of maps on vertices, edges and labels,
commuting with all three maps. A marked system consists of two graph morphisms
\(\mathcal T\to\mathcal L\to\mathcal X\). Its morphisms consist of three graph
morphisms commuting with both marking squares.
\end{definition}

% These are actual category laws, not a claimed equivalence with an unformalized functor category.
\begin{theorem}[The two category structures]
Labelled graphs and their morphisms form a category. Marked systems and their
square-preserving morphisms form a category. Neither construction silently
requires injective marks, determinism, invariant targets, order, or measure.
\end{theorem}
\begin{proof}
The graph identity is identity on vertices, edges and labels. Composition is
composition of each of these three functions. For example, composing the source
square for \(f\) with that for \(g\) gives
\(g_Vf_Vs=s g_Ef_E\); the target and label squares are identical calculations.
Function associativity and identity give the graph category laws.
For marked systems take identities and composition at all three stages.
If \(f\) and \(g\) preserve a marking map \(j\), associativity and their two
squares give \(j;f;g=f;j';g=f;g;j''\). Apply this argument to each mark.
The component function laws then give associativity and both identity laws,
including the two compatibility squares. These calculations use none of the
additional properties listed in the statement.
\end{proof}

\section{Target reflection and variable clocks}

% A commuting map does not automatically reflect a marked target.
\begin{theorem}[Target-reflecting semiconjugacy]
Suppose \(h\circ f=g\circ h\), and predicates \(A,B\) satisfy
\(B(h(x))\Rightarrow A(x)\) for every \(x\).
Then \(h(f^n(x))=g^n(h(x))\) for every \(n\ge0\). If the latter orbit ever enters
\(B\), the original orbit enters \(A\) at the same step.
\end{theorem}
\begin{proof}
The iterate identity follows from the zero-step identity and induction, applying
the commuting equation at each successor. If \(B(g^n(h(x)))\), replace its
argument by \(h(f^n(x))\) and apply target reflection at that one source state.
This uses a single common capture time. No injectivity or surjectivity is assumed.
\end{proof}

% Positive clocks are needed for unbounded raw time, but not for the finite forward-capture identity.
\begin{theorem}[Clock]
Let the clock \(\tau\) be natural-valued and suppose
\[
h(f(x))=g^{\tau(x)}(h(x)).
\]
Define
\[
T_0(x)=0,\qquad T_{n+1}(x)=T_n(x)+\tau(f^n(x)).
\]
Then \(h(f^n(x))=g^{T_n(x)}(h(x))\).
If \(A(x)\Rightarrow B(h(x))\), capture of the \(f\)-orbit in \(A\) implies
capture of the \(g\)-orbit in \(B\). If \(\tau(x)\ge1\) at every state, then
\(T_n(x)\ge n\). No reverse capture assertion is made.
\end{theorem}
\begin{proof}
Induction on \(n\), using \(g^{a+b}=g^b\circ g^a\), gives the exact iterate
identity. A source capture at time \(n\) maps into \(B\) and occurs in the raw
orbit at time \(T_n(x)\). Positivity gives
\(T_{n+1}(x)\ge T_n(x)+1\); induction from \(T_0=0\) proves the lower bound.
\end{proof}

\section{Counterexamples}

% Signed division invalidates an unsigned absorption argument.
\begin{theorem}[A nonzero signed cycle]
For \((q,c,d)=(2,2,-2)\), the orbit of \(1\) is the cycle
\(1\to4\to-2\to1\). In particular, this orbit does not reach zero.
\end{theorem}
\begin{proof}
Since \(-2\nmid1\), the first image is \(2\cdot1+2=4\).
The next two images are exact divisions: \(4/(-2)=-2\) and
\((-2)/(-2)=1\). Thus the set \(\{1,4,-2\}\) is forward-invariant.
Induction on the number of steps places every iterate in this set, which excludes zero.
\end{proof}

% Keeping c is necessary even for a one-letter exponent word.
\begin{theorem}[Offset dependence]
For \((q,c,d)=(3,5,2)\), the raw word \(ADD\) closes at \(5\).
Its exponent-two presentation fixes \(5\); changing \(c\) to one does not.
\end{theorem}
\begin{proof}
The legal raw steps are \(5\to20\to10\to5\): the first state is odd and
both division states are even. The rational presentation is
\((3x+5)/4\), so its value at \(5\) is \(20/4=5\).
After replacing the offset by one, the value is \(16/4=4\ne5\).
\end{proof}

% Scalar closure leaves the guards of zero-exponent words unresolved.
\begin{theorem}[An illegal rational fixed word]
For \((q,c,d)=(3,1,2)\), the exponent word \((0,0,3,4)\) fixes \(1\)
rationally, but its raw expansion is not legal at \(1\).
\end{theorem}
\begin{proof}
Its data are \(P=3^4=81\), \(Q=2^7=128\), and
\(B=27+9+3+8=47\). Hence \((81\cdot1+47)/128=1\).
The raw expansion starts with two affine letters. Its first letter sends
\(1\) to \(4\), where the division branch is mandatory. The second affine
letter is therefore illegal.
\end{proof}

% Word-specific noncoprime scaling can work without defining a global conjugacy.
\begin{theorem}[Scaling]
For parameters
\[
(q,c,d)=(-5,31,-6),
\]
the word \(AAD\) closes at \(4\).
Multiplying the offset and state by two preserves this word;
multiplying both by three makes its first letter illegal.
\end{theorem}
\begin{proof}
The unscaled steps are \(4\to11\to-24\to4\).
Neither \(4\) nor \(11\) is divisible by six, and \(-24\) is divisible by \(-6\).
At offset \(62\), the steps are \(8\to22\to-48\to8\), with the same legal guards.
At offset \(93\), the proposed starting state is \(12\); because \(-6\mid12\),
its first instruction must be a division, not the proposed affine letter.
\end{proof}

\bibliographystyle{amsalpha}
\bibliography{Foundations}

\end{document}
